\documentclass[11pt, a4paper]{article}

% Gemeinsame Style-Definitionen (TEXINPUTS muss templates/ enthalten — siehe scripts/build_tex.py)
% bfw-style.tex — gemeinsame Style-Definitionen für alle Handouts/Aufgabenhefte
% Voraussetzung: Kompilation mit xelatex (wegen fontspec)

\usepackage[ngerman]{babel}
\usepackage{geometry}
\geometry{a4paper, top=2.2cm, bottom=2.2cm, left=2.3cm, right=2.3cm}

\usepackage{fontspec}
% Fallback-Kette: Source Sans Pro -> Calibri -> Segoe UI -> Latin Modern Sans
% (Latin Modern ist immer mit MiKTeX vorhanden)
\IfFontExistsTF{Source Sans Pro}{%
  \setmainfont{Source Sans Pro}[Ligatures=TeX]%
}{\IfFontExistsTF{Calibri}{%
  \setmainfont{Calibri}[Ligatures=TeX]%
}{\IfFontExistsTF{Segoe UI}{%
  \setmainfont{Segoe UI}[Ligatures=TeX]%
}{%
  \setmainfont{Latin Modern Sans}[Ligatures=TeX]%
}}}
\IfFontExistsTF{Source Code Pro}{%
  \setmonofont{Source Code Pro}[Scale=0.92]%
}{\IfFontExistsTF{Consolas}{%
  \setmonofont{Consolas}[Scale=0.92]%
}{%
  \setmonofont{Latin Modern Mono}[Scale=0.92]%
}}

\usepackage{xcolor}
% Farbpalette: hell, freundlich, modern
\definecolor{bfwPrimary}{HTML}{0EA5A4}    % Petrol/Türkis - Primär
\definecolor{bfwPrimaryDark}{HTML}{0F766E} % dunkler Petrol für Headings
\definecolor{bfwAccent}{HTML}{F59E0B}     % Amber - Akzent
\definecolor{bfwSuccess}{HTML}{10B981}    % Grün - Erfolg / Lösung
\definecolor{bfwWarn}{HTML}{F97316}       % Orange - Achtung
\definecolor{bfwInk}{HTML}{0F172A}        % fast Schwarz
\definecolor{bfwInkSoft}{HTML}{475569}    % gedämpftes Grau
\definecolor{bfwBgSoft}{HTML}{F8FAFC}     % off-white
\definecolor{bfwBgTeal}{HTML}{ECFEFF}     % helles Türkis
\definecolor{bfwBgAmber}{HTML}{FEF3C7}    % helles Gelb
\definecolor{bfwBgRose}{HTML}{FFE4E6}     % helles Rosé für Eselsbrücken

\usepackage[colorlinks=true, linkcolor=bfwPrimaryDark, urlcolor=bfwPrimaryDark]{hyperref}
\usepackage{enumitem}
\setlist[itemize]{leftmargin=*, itemsep=2pt, topsep=2pt}
\setlist[enumerate]{leftmargin=*, itemsep=2pt, topsep=2pt}

\usepackage[most]{tcolorbox}
\usepackage{booktabs}
\usepackage{tikz}
\usetikzlibrary{decorations.pathmorphing, positioning, shapes.geometric, arrows.meta}

% Merkkästchen — wichtige Definition / Kernpunkt
\newtcolorbox{merksatz}[1][]{%
  enhanced, breakable,
  colback=bfwBgTeal,
  colframe=bfwPrimary,
  boxrule=0.6pt,
  arc=4pt,
  left=10pt, right=10pt, top=8pt, bottom=8pt,
  fonttitle=\bfseries\color{bfwPrimaryDark},
  title={\faLightbulb\ Merksatz},
  #1
}

% Eselsbrücke — Mnemoniks
\newtcolorbox{eselsbruecke}[1][]{%
  enhanced, breakable,
  colback=bfwBgRose,
  colframe=bfwAccent,
  boxrule=0.6pt,
  arc=4pt,
  left=10pt, right=10pt, top=8pt, bottom=8pt,
  fonttitle=\bfseries\color{bfwWarn},
  title={Eselsbrücke},
  #1
}

% Beispiel-Box
\newtcolorbox{beispiel}[1][]{%
  enhanced, breakable,
  colback=bfwBgSoft,
  colframe=bfwInkSoft,
  boxrule=0.4pt,
  arc=3pt,
  left=10pt, right=10pt, top=8pt, bottom=8pt,
  fonttitle=\bfseries\color{bfwInk},
  title={Beispiel},
  #1
}

% Lösungs-Box (für Aufgabenhefte)
\newtcolorbox{loesung}[1][]{%
  enhanced, breakable,
  colback=bfwBgAmber!40,
  colframe=bfwSuccess,
  boxrule=0.6pt,
  arc=3pt,
  left=10pt, right=10pt, top=8pt, bottom=8pt,
  fonttitle=\bfseries\color{bfwSuccess!70!black},
  title={Lösung},
  #1
}

% Glossar-Eintrag
\newcommand{\glossareintrag}[2]{%
  \par\noindent\textbf{\color{bfwPrimaryDark}#1} \enspace #2\par\smallskip
}

% Section-Stil — etwas farbiger als Standard
\usepackage{titlesec}
\titleformat{\section}
  {\Large\bfseries\color{bfwPrimaryDark}}
  {\thesection}{0.6em}{}
\titleformat{\subsection}
  {\large\bfseries\color{bfwPrimary}}
  {\thesubsection}{0.5em}{}
\titleformat{\subsubsection}
  {\normalsize\bfseries\color{bfwInk}}
  {\thesubsubsection}{0.4em}{}

% TikZ-Stil "skizzenhaft" — etwas weicher als Rough.js, aber stabil
\tikzset{
  roughbox/.style = {
    draw=bfwPrimaryDark, thick, fill=bfwBgTeal,
    rounded corners=4pt, inner sep=6pt
  },
  roughaccent/.style = {
    draw=bfwAccent, thick, fill=bfwBgAmber,
    rounded corners=4pt, inner sep=6pt
  },
  roughline/.style = {
    draw=bfwInkSoft, thick, -{Stealth[length=2.5mm]}
  }
}

% Bullet-Symbole (bewusst kein FontAwesome — vermeidet Font-Installations-Probleme)
\newcommand{\faLightbulb}{$\bullet$}
\newcommand{\faBookmark}{$\bullet$}

% Header/Footer
\usepackage{fancyhdr}
\pagestyle{fancy}
\fancyhf{}
\renewcommand{\headrulewidth}{0pt}
\fancyfoot[L]{\small\color{bfwInkSoft}\thepage}
\fancyfoot[R]{\small\color{bfwInkSoft} BFW M-064 Beschaffung im Büromanagement}


\title{\vspace*{-1.5cm}\Huge\bfseries\color{bfwPrimaryDark}Tag~9: Kommunikationsformen \& Teamarbeit\\[0.4em]
  \large\color{bfwInkSoft}\textnormal{Wie wir verstanden werden --- und gut zusammenarbeiten}}
\author{\textsf{Büroprozesse gestalten und Arbeitsvorgänge organisieren \textperiodcentered{} M-067}\\
\textsf{\small\copyright\ Can Siebert\,\textperiodcentered\,2026}}
\date{09.07.2026}

\begin{document}
\maketitle
\thispagestyle{empty}

\vspace{1em}
\begin{merksatz}[title={\bfseries Worum geht es heute?}]
Kein Büroprozess funktioniert ohne Kommunikation --- und kaum eine Aufgabe wird heute
allein erledigt. Wer eine Nachricht sendet, will verstanden werden; doch zwischen
„gemeint” und „angekommen” liegt viel Raum für Missverständnisse. In diesem Handout
lernen Sie, wie Kommunikation funktioniert (Sender-Empfänger-Modell, die vier Seiten
einer Nachricht, die Axiome von Watzlawick), welche Kommunikations\-formen das Büro
kennt --- vom persönlichen Gespräch über das Telefon bis zur E-Mail und Video\-konferenz ---
und wie aus einer Gruppe ein leistungsfähiges Team wird. Das ist Grundlage für jede
Zusammenarbeit und für die IHK-Abschlussprüfung.
\end{merksatz}

\tableofcontents
\vspace{1em}
\hrule
\vspace{1em}

\section{Wie Kommunikation funktioniert}

Kommunikation ist der Austausch von Informationen zwischen mindestens zwei Personen.
Im Büro ist sie das wichtigste Werkzeug überhaupt: Aufträge, Absprachen, Auskünfte und
Entscheidungen entstehen durch Kommunikation. Damit sie gelingt, lohnt sich ein Blick
darauf, was beim Senden und Empfangen tatsächlich passiert.

\subsection{Das Sender-Empfänger-Modell}

Das \textbf{Sender-Empfänger-Modell} beschreibt Kommunikation als Übertragung einer
Nachricht. Der \emph{Sender} hat eine Absicht und \emph{verschlüsselt} sie in Zeichen
(Worte, Tonfall, Körpersprache) --- das nennt man \emph{Encodieren}. Diese Nachricht
wird über einen \emph{Kanal} (Luft, Telefon, E-Mail) übertragen. Der \emph{Empfänger}
nimmt sie auf und \emph{entschlüsselt} sie zu einer eigenen Bedeutung
(\emph{Decodieren}). Anschließend gibt er idealerweise eine Rückmeldung
(\emph{Feedback}), sodass der Sender prüfen kann, ob er verstanden wurde.

\begin{merksatz}
Zwischen „gemeint”, „gesagt”, „gehört” und „verstanden” liegen vier verschiedene Dinge.
Jede Störung auf dem Weg --- man nennt sie \emph{Rauschen} --- kann dazu führen, dass die
Nachricht beim Empfänger anders ankommt, als sie gemeint war.
\end{merksatz}

\emph{Rauschen} (Störungen) entsteht z.\,B.\ durch Lärm im Großraumbüro, eine schlechte
Telefonverbindung, Ablenkung, unklare Formulierungen oder durch unterschiedliches
Vorwissen, Erfahrungen und Stimmungen von Sender und Empfänger. Weil Encodieren und
Decodieren immer \emph{subjektiv} sind, sind Missverständnisse normal --- und gutes
Nachfragen ist keine Schwäche, sondern Qualitätssicherung.

\subsection{Die vier Seiten einer Nachricht}

Der Psychologe Friedemann Schulz von Thun hat gezeigt, dass jede Nachricht
\emph{vier Seiten} hat --- der Sender spricht mit „vier Schnäbeln”, der Empfänger hört
mit „vier Ohren”:

\begin{enumerate}
  \item \textbf{Sachebene} --- die reine Information, der sachliche Inhalt (\emph{Worüber}
    informiere ich?).
  \item \textbf{Selbstoffenbarung} --- was der Sender von sich preisgibt (\emph{Was}
    \emph{sage ich über mich?}).
  \item \textbf{Beziehungsebene} --- was der Sender vom Empfänger hält und wie er zu ihm
    steht (\emph{Wie stehen wir zueinander?}).
  \item \textbf{Appellebene} --- wozu der Sender den Empfänger veranlassen will
    (\emph{Wozu möchte ich dich bringen?}).
\end{enumerate}

\begin{beispiel}
Eine Kollegin sagt: \emph{„Da ist ein Fehler in der Summenzeile.”}
\begin{itemize}[itemsep=0pt]
  \item \emph{Sache:} In der Summenzeile ist ein Rechenfehler.
  \item \emph{Selbstoffenbarung:} „Mir ist Genauigkeit wichtig, mir ist es aufgefallen.”
  \item \emph{Beziehung:} „Ich schaue auf deine Arbeit” --- kann als „Ich traue dir keine
    fehlerfreie Arbeit zu” gehört werden.
  \item \emph{Appell:} „Bitte korrigiere die Summenzeile.”
\end{itemize}
Hört der Kollege mit dem \emph{Beziehungsohr}, fühlt er sich kritisiert --- und es
entsteht ein Konflikt, obwohl es sachlich nur um eine Zahl ging.
\end{beispiel}

\subsection{Die fünf Axiome von Watzlawick}

Paul Watzlawick hat fünf Grundregeln (\emph{Axiome}) der Kommunikation formuliert. Die
wichtigsten für den Büroalltag:

\begin{enumerate}
  \item \textbf{Man kann nicht \emph{nicht} kommunizieren.} Auch Schweigen, ein
    Gesichtsausdruck oder eine nicht beantwortete E-Mail sind Botschaften.
  \item \textbf{Jede Kommunikation hat einen Inhalts- und einen Beziehungsaspekt}, wobei
    die Beziehung den Inhalt bestimmt (vgl.\ die vier Seiten einer Nachricht).
  \item \textbf{Kommunikation ist Ursache und Wirkung zugleich} (\emph{Interpunktion}):
    Jeder sieht sich als Reagierender --- „Ich bin nur so kurz angebunden, weil du nie
    antwortest” gegen „Ich antworte nicht, weil du so kurz angebunden bist”.
  \item \textbf{Digitale und analoge Kommunikation:} Inhalte werden digital (über Worte)
    und analog (über Körpersprache, Tonfall) übertragen.
  \item \textbf{Kommunikation ist symmetrisch oder komplementär}, also auf Augenhöhe oder
    durch ein Gefälle geprägt (z.\,B.\ Vorgesetzte/Mitarbeitende).
\end{enumerate}

\subsection{Verbal, nonverbal und paraverbal}

Wir kommunizieren nicht nur mit Worten. Drei Ebenen wirken immer zusammen:

\begin{center}
\begin{tabular}{p{3.2cm} p{4.2cm} p{5.2cm}}
\toprule
\textbf{Ebene} & \textbf{Was?} & \textbf{Beispiele}\\
\midrule
Verbal & der Inhalt, die Worte & Wortwahl, Satzbau, das Gesagte\\
Nonverbal & Körpersprache ohne Worte & Mimik, Gestik, Körperhaltung, Blickkontakt\\
Paraverbal & die \emph{Art} des Sprechens & Tonfall, Lautstärke, Sprechtempo, Betonung, Pausen\\
\bottomrule
\end{tabular}
\end{center}

\medskip
Passen die Ebenen nicht zusammen (\emph{Inkongruenz}) --- etwa „Mir geht es gut” mit
gesenktem Blick und leiser Stimme ---, entsteht eine widersprüchliche Doppelbotschaft.
Empfänger glauben im Zweifel der nonverbalen und paraverbalen Ebene; die Glaubwürdigkeit
leidet.

\begin{eselsbruecke}
\textbf{„Was, Wie, Womit”} --- \emph{verbal} = \emph{Was} ich sage, \emph{paraverbal} =
\emph{Wie} ich es sage (Stimme), \emph{nonverbal} = \emph{Womit} mein Körper es begleitet
(Mimik, Gestik). Studien betonen: Der Ton macht oft mehr als das Wort.
\end{eselsbruecke}

\section{Kommunikationsformen im Büro}

Für jeden Anlass gibt es eine passende Kommunikationsform. Die Kunst besteht darin, sie
\emph{situationsgerecht} auszuwählen --- nach Vertraulichkeit, Dokumentationsbedarf,
Geschwindigkeit und Zahl der Beteiligten.

\subsection{Mündlich und schriftlich}

\textbf{Mündliche Kommunikation} (persönliches Gespräch, Besprechung) ist schnell,
ermöglicht sofortiges Nachfragen und überträgt auch Tonfall und Körpersprache --- ideal
für Vertrauliches, Kritik und komplexe Themen. Nachteil: Es bleibt kein Nachweis.
\textbf{Schriftliche Kommunikation} (Brief, E-Mail) ist dagegen \emph{dokumentiert und
beweissicher}, lässt sich in Ruhe formulieren und an viele Empfänger verteilen ---
dafür fehlen Tonfall und Körpersprache, und es kann zu Verzögerungen kommen.

\subsection{Telefon und Telefonnotiz}

Das Telefon verbindet Schnelligkeit mit der persönlichen Note der Stimme. Weil das
Gespräch flüchtig ist, ist die \textbf{Telefonnotiz} wichtig: Wer hat \emph{wann} für
\emph{wen} angerufen, \emph{worum} ging es, \emph{was} ist zu tun (Rückruf?), und
\emph{wer} hat die Notiz aufgenommen? So gehen Informationen und Aufträge nicht verloren.

\subsection{Elektronische Post --- die geschäftliche E-Mail}

In vielen Unternehmen erfolgt die Korrespondenz auf elektronischem Wege mittels E-Mails,
die den postalischen Geschäftsbrief ersetzen. Aus diesem Grund müssen geschäftliche
E-Mails bestimmte \textbf{Pflichtangaben} enthalten. Dies ist im \emph{Gesetz über das
elektronische Handelsregister und Genossenschaftsregister sowie das Unternehmens\-register
(EHUG)} geregelt: Der Empfänger geschäftlicher E-Mails soll über den Absender im gleichen
Umfang informiert werden wie bei einem postalischen Geschäftsbrief.

\begin{merksatz}
Auf allen Geschäftsbriefen --- dazu zählen auch Bestellungen, Lieferscheine, Rechnungen
und Quittungen in Form von E-Mail oder Telefax --- sind bestimmte Angaben zum Unternehmen
aufzuführen. \emph{Welche} Pflichtangaben eine E-Mail enthalten muss, hängt von der
jeweiligen \emph{Rechtsform} des Unternehmens ab.
\end{merksatz}

Typische Pflichtangaben sind --- je nach Rechtsform --- z.\,B.\ Firma, Rechtsform,
Sitz der Gesellschaft, Registergericht und Handelsregisternummer sowie die Geschäfts\-führer
bzw.\ Vorstandsmitglieder. Die genaue Adresse, Telefon- oder Faxnummern, ein Postfach
oder die Bankverbindung zählen \emph{nicht} zu den Pflichtangaben in einem Geschäftsbrief.
Wird ein E-Mail-Vordruck auch als Rechnung verwendet, sind zusätzlich die Pflichtangaben
nach \S~14 Umsatzsteuergesetz (UStG) zu beachten (z.\,B.\ die
Umsatzsteueridentifikations\-nummer).

\subsubsection*{E-Mail-Knigge --- die wichtigsten Regeln}

\begin{itemize}
  \item Aussagekräftiger \textbf{Betreff} --- der Empfänger erkennt sofort, worum es geht.
  \item Passende \textbf{Anrede} und Grußformel, höflicher und sachlicher Ton.
  \item \textbf{Kurz und klar} formulieren; eine E-Mail --- ein Anliegen.
  \item Empfänger bewusst wählen: \emph{An} nur an Verantwortliche, \emph{CC} sparsam,
    \emph{BCC} bei Verteilern.
  \item Rechtschreibung prüfen; \textbf{Anhänge} ankündigen und tatsächlich beifügen.
  \item \textbf{Signatur} mit den Pflichtangaben nicht vergessen.
  \item Zeitnah antworten, aber nicht überstürzt; bei Ärger erst entwerfen, dann senden.
\end{itemize}

\subsection{Messenger/Chat und Videokonferenz}

\textbf{Messenger und Chat} (z.\,B.\ Teams, Slack) eignen sich für kurze, informelle und
schnelle Abstimmungen im Team; sie ersetzen aber keine verbindliche, dokumentierte
Korrespondenz. \textbf{Videokonferenzen} verbinden mehrere Standorte, sparen Reisezeit
und übertragen Bild und Ton --- ideal für verteilte Teams. Wichtig sind klare Regeln
(Mikrofon stummschalten, Wortmeldungen, Agenda, Protokoll), damit die Besprechung
strukturiert bleibt.

\section{Teamarbeit}

\emph{Allein ist man schneller, gemeinsam kommt man weiter.} Ein \textbf{Team} ist mehr
als eine Gruppe: Seine Mitglieder verfolgen ein \emph{gemeinsames Ziel}, ergänzen sich
mit unterschiedlichen Fähigkeiten und Rollen und fühlen sich \emph{gemeinsam}
verantwortlich.

\subsection{Vor- und Nachteile der Teamarbeit}

\begin{center}
\begin{tabular}{p{6.2cm} p{6.2cm}}
\toprule
\textbf{Vorteile} & \textbf{Nachteile}\\
\midrule
Mehr Ideen und Kreativität & Höherer Abstimmungs- und Zeitaufwand\\
Stärken ergänzen sich, Arbeitsteilung & Konfliktpotenzial zwischen Mitgliedern\\
Höhere Akzeptanz von Entscheidungen & Gefahr des sozialen Faulenzens\\
Motivation, gegenseitige Vertretung & Gruppendruck / Gruppendenken\\
\bottomrule
\end{tabular}
\end{center}

\medskip
\textbf{Soziales Faulenzen} (Trittbrettfahren) bedeutet, dass Einzelne ihre Leistung
zurückfahren, weil sie sich auf das Team verlassen. Dagegen helfen klare Rollen,
messbare Einzelbeiträge, eine überschaubare Teamgröße und transparente Ziele.

\subsection{Teamrollen}

Erfolgreiche Teams sind \emph{gemischt}: Sie brauchen Ideengeber, Umsetzer, Koordinatoren,
kritische Prüfer und Personen, die auf das Miteinander achten. Wenn alle dieselbe Rolle
einnehmen wollen, entstehen entweder Leerstellen oder Reibung. Vielfalt in Persönlichkeit,
Erfahrung und Stärken ist daher ein Vorteil --- vorausgesetzt, die Rollen sind geklärt.

\subsection{Teamphasen nach Tuckman}

Teams entwickeln sich in typischen Phasen. Bruce Tuckman beschreibt fünf:

\begin{enumerate}
  \item \textbf{Forming} (Orientierung) --- das Team findet sich; höflich, unsicher,
    Ziele und Rollen sind noch unklar.
  \item \textbf{Storming} (Konflikt) --- Auseinandersetzungen über Rollen, Aufgaben und
    Einfluss; Reibung ist normal und notwendig.
  \item \textbf{Norming} (Organisation) --- gemeinsame Regeln und Rollen werden gefunden,
    der Zusammenhalt wächst.
  \item \textbf{Performing} (Leistung) --- das Team arbeitet effizient, selbstorganisiert
    und zielgerichtet.
  \item \textbf{Adjourning} (Auflösung) --- Abschluss, Auswertung und Auflösung des Teams
    (z.\,B.\ am Projektende).
\end{enumerate}

\begin{eselsbruecke}
\textbf{„Forming, Storming, Norming, Performing, Adjourning”} --- die Reihenfolge merkt
man sich am Reim auf \emph{-orming}. Wichtig: Die \emph{Storming}-Phase ist kein
Scheitern, sondern ein notwendiger Schritt auf dem Weg zur Leistungsfähigkeit.
\end{eselsbruecke}

\subsection{Konflikte und Feedbackregeln}

Konflikte gehören zu jedem Team. Entscheidend ist, sie \emph{früh und sachlich}
anzusprechen, statt sie zu vermeiden. Gutes \textbf{Feedback} hilft dabei:

\begin{itemize}
  \item \textbf{Zeitnah und konkret} --- auf beobachtbares Verhalten bezogen, nicht
    pauschal („immer”, „nie”).
  \item \textbf{Ich-Botschaften} statt Du-Vorwürfe: „Ich habe den Eindruck \dots” statt
    „Du machst immer \dots”.
  \item \textbf{Beschreibend statt bewertend} --- Verhalten schildern, nicht die Person
    abwerten.
  \item \textbf{Wertschätzend und konstruktiv} --- mit einem konkreten
    Veränderungsvorschlag.
  \item Der \emph{Empfänger} lässt ausreden, rechtfertigt sich nicht sofort und fragt bei
    Unklarheit nach.
\end{itemize}

\section{Zusammenfassung}

\begin{enumerate}
  \item Das \textbf{Sender-Empfänger-Modell} erklärt Kommunikation als Encodieren und
    Decodieren --- \emph{Rauschen} verursacht Missverständnisse.
  \item Jede Nachricht hat \textbf{vier Seiten}: Sache, Selbstoffenbarung, Beziehung,
    Appell --- gehört mit „vier Ohren”.
  \item Watzlawick: \textbf{Man kann nicht nicht kommunizieren}; jede Kommunikation hat
    einen Inhalts- und einen Beziehungsaspekt.
  \item \textbf{Verbal, nonverbal, paraverbal} wirken zusammen --- bei Widerspruch zählt
    der Eindruck von Körpersprache und Stimme.
  \item \textbf{Kommunikationsformen} im Büro situationsgerecht wählen; geschäftliche
    \textbf{E-Mails} brauchen Pflichtangaben (EHUG) und einen sauberen E-Mail-Knigge.
  \item \textbf{Teamarbeit} bietet Vorteile (Ideen, Arbeitsteilung) und Risiken (Konflikte,
    Faulenzen); Teams durchlaufen die \textbf{Tuckman-Phasen} Forming bis Adjourning.
  \item \textbf{Feedbackregeln} (Ich-Botschaften, konkret, wertschätzend) machen Konflikte
    lösbar.
\end{enumerate}

\newpage

\section*{Glossar}
\addcontentsline{toc}{section}{Glossar}

\glossareintrag{Kommunikation}{Austausch von Informationen zwischen mindestens zwei Personen über einen Kanal.}

\glossareintrag{Sender-Empfänger-Modell}{Grundmodell der Kommunikation: Der Sender encodiert eine Nachricht, der Empfänger decodiert sie und gibt Feedback.}

\glossareintrag{Encodieren / Decodieren}{Verschlüsseln einer Absicht in Zeichen (Sender) bzw.\ Entschlüsseln der Zeichen zu einer Bedeutung (Empfänger).}

\glossareintrag{Rauschen}{Jede Störung, die die Übertragung einer Nachricht beeinträchtigt (Lärm, Ablenkung, unklare Worte, Vorurteile).}

\glossareintrag{Vier Seiten einer Nachricht}{Modell von Schulz von Thun: Jede Nachricht hat eine Sach-, Selbstoffenbarungs-, Beziehungs- und Appellebene.}

\glossareintrag{Watzlawick-Axiome}{Fünf Grundregeln der Kommunikation, u.\,a.\ „Man kann nicht nicht kommunizieren”.}

\glossareintrag{Verbale Kommunikation}{Der inhaltliche Teil: die gesprochenen oder geschriebenen Worte.}

\glossareintrag{Nonverbale Kommunikation}{Körpersprache ohne Worte: Mimik, Gestik, Körperhaltung, Blickkontakt.}

\glossareintrag{Paraverbale Kommunikation}{Die Art des Sprechens: Tonfall, Lautstärke, Sprechtempo, Betonung, Pausen.}

\glossareintrag{Telefonnotiz}{Schriftliche Festhaltung eines Telefonats (wer, wann, für wen, worum, was ist zu tun).}

\glossareintrag{EHUG}{Gesetz über das elektronische Handelsregister; regelt u.\,a.\ Pflichtangaben in geschäftlichen E-Mails.}

\glossareintrag{Pflichtangaben}{Vorgeschriebene Angaben zum Unternehmen auf Geschäftsbriefen und geschäftlichen E-Mails (abhängig von der Rechtsform).}

\glossareintrag{E-Mail-Knigge}{Regeln für höfliche, klare und korrekte geschäftliche E-Mails (Betreff, Anrede, Anhänge, Signatur).}

\glossareintrag{Team}{Personengruppe mit gemeinsamem Ziel, ergänzenden Rollen und gemeinsamer Verantwortung.}

\glossareintrag{Soziales Faulenzen}{Leistungsrückgang Einzelner, weil sie sich auf das Team verlassen (Trittbrettfahren).}

\glossareintrag{Teamphasen nach Tuckman}{Forming, Storming, Norming, Performing, Adjourning --- die typische Entwicklung eines Teams.}

\glossareintrag{Feedback}{Rückmeldung über beobachtetes Verhalten; konstruktiv, konkret, zeitnah und in Ich-Botschaften.}

\end{document}
